make a minimal symbolic agent framework for studying agents across many environments under one abstraction.

the core idea is:

an agent receives symbols from an environment, updates internal symbolic state, and emits symbols back to the environment.

define the interaction approximately as:

observation symbols → agent → action symbols → environment → reward + next observation symbols

the framework should be general enough that many apparently different learning/reasoning paradigms can be expressed as instances of the same agent-environment loop.

for example:

* language modeling:

  * observation: previous tokens
  * action: predicted next-token
  * environment transition: append/advance to next token
  * reward: log probability / negative log-likelihood-derived reward

* images can be pre-processed by SAM and/or other engines to make a list of bbox+class symbols for the image. this still requires learning how to integrate symbolic elements of the lowest layer of the scene graph into a integrated scene understanding

* classification:

  * observation: symbolic representation of an example
  * action: class symbol
  * reward: correctness

* theorem proving:

  * observation: axioms, current proof state, goal
  * action: inference-rule application / generated proposition
  * reward: proof completion and/or neg log duration

* program synthesis:

  * observation: specification + current program
  * action: symbolic program edit / ast operation
  * reward: test results

* planning:

  * observation: symbolic world state
  * action: symbolic world action
  * reward: task utility

* games:

  * observation: symbolic game state
  * action: symbolic move
  * reward: game reward

* tool-using agents:

  * observation: messages (containing symbols for their language contents) + tool results (containing multiple symbols for their fn class + params)
  * action: symbolic tool call or response
  * reward: task grader

do NOT build a huge ontology, theorem proving platform, knowledge graph system, or 'general' symbolic-ai architecture that doesn't do anything.

focus specifically on the agent/environment abstraction.

the most important interfaces should be small.

something conceptually like:

```python
class Symbol:
    type: str
    value: Any
    metadata: dict

class Environment:
    reset() -> Observation
    step(action: Action) -> Transition

class Transition:
    observation: Observation
    reward: float
    terminated: bool
    info: dict

class Agent:
    reset()
    act(observation: Observation) -> Action
```

but do improve these interfaces where necessary.

observations and actions should fundamentally be symbolic structures, not assumed to be tensors.

perhaps a symbol may be:

* token
* string
* number
* identifier
* predicate
* relation
* ast node
* tuple
* list
* tree
* graph reference
* function/operator application
* structured record

support composition so complex observations/actions can be constructed from simpler symbols.

there should be an interchangeable encoding layer:

raw environment state
→ encoder / lexer
→ symbols
→ agent

and:

agent symbols
→ decoder / interpreter
→ environment action

this lets different environments expose radically different underlying state while presenting a common symbolic interface to the agent.

likewise, agents should be interchangeable.

start with several deliberately different agent implementations:

1. random symbolic policy
2. lookup/table policy
3. rule-based agent
4. finite-state agent
5. search/planning agent
6. simple learned neural policy over symbolic inputs
7. autoregressive token policy

the point is not to force them to share the same internal architecture. they only need to satisfy the same external symbolic interaction contract.

keep internal state explicit:

agent state_t + observation_t
→ agent state_{t+1} + action_t

ideally model the agent mathematically as:

A : (S_agent, O*) → (S_agent, A*)

where O* and A* are structured symbolic observations and actions.

also make environments composable around the same contract.

implement a handful of small environments showing that the abstraction really unifies different paradigms:

* next-token prediction
* sequence completion
* classification
* boolean/logical inference
* simple gridworld with symbolic observations/actions
* arithmetic-expression completion
* simple theorem/proof environment
* simple tool-call environment

for next-token prediction, explicitly demonstrate how standard language modeling fits the environment abstraction.

for a sequence x_1 ... x_n:

observation_t = x_1 ... x_t
action_t = prediction for x_{t+1}

and define reward from the predicted probability assigned to the true next symbol, e.g.:

r_t = log p(x_{t+1} | x_{\le t})

or equivalently treat negative log-likelihood as cost.

this example matters because one research objective is to show that supervised prediction, reinforcement learning, reasoning, planning, and agentic interaction can often be represented using the same sequential agent/environment formalism.

separate three concepts carefully:

* symbol space
* agent policy
* environment dynamics/reward

do not bake task semantics into the agent.

make symbol spaces typed and inspectable. each environment should declare its observation-symbol space and action-symbol space.

support both:

* discrete action symbols
* structured action symbols

for example:

`LEFT`

and:

`Call(function="search", args={"query": "..."})`

might both be valid actions.

build trajectories as first-class objects:

trajectory = [
(observation_0, action_0, reward_0),
(observation_1, action_1, reward_1),
...
]

make them serializable and easy to inspect.

provide a minimal training/evaluation abstraction, but DO NOT prematurely commit the architecture to one optimizer.

we may later want to train agents using:

* maximum likelihood
* behavioral cloning
* policy gradients
* q-learning
* search
* evolutionary methods
* symbolic rule induction
* offline rl
* online rl

the environment contract should not care.

include a small experiment harness where i can run something like:

```bash
symbolic-agent run \
  --env next_token \
  --agent autoregressive
```

or:

```bash
symbolic-agent run \
  --env logic \
  --agent rule_based
```

and compare combinations.

the architectural emphasis is:

environment semantics are interchangeable.
symbol encoders/decoders are interchangeable.
agent implementations are interchangeable.

the invariant is only the symbolic interaction protocol.

keep the implementation small, legible, and hackable. prefer a few hundred lines of strong abstractions and working examples over thousands of lines of enterprise scaffolding.

before coding, briefly inspect existing agent-environment abstractions such as gymnasium and modern agentic-rl environment frameworks, including prime intellect's verifiers/research-environments approach, and borrow useful lifecycle ideas without inheriting unnecessary complexity.

prime intellect's useful pattern is that heterogeneous tasks can retain their native setup and grading semantics while being normalized behind a common contract. apply the analogous principle here: environments should preserve their own transition and reward semantics while exposing them through one symbolic-agent interface.

the resulting project should make this question easy to investigate:

"how much of ai can be represented as an agent consuming symbols, producing symbols, and receiving feedback from an environment?"

addendum: make the system trainable and evaluable against REAL existing datasets and environment zoos.

do not stop once the symbolic agent/environment abstraction works on toy environments. immediately use it to wrap substantial existing training and evaluation distributions.

first, build a generic dataset-backed environment adapter for language modeling and other static supervised datasets.

the important abstraction is that a dataset becomes an environment by defining:

dataset sample
→ initial symbolic observation
→ one or more symbolic actions
→ deterministic/stochastic transition
→ reward/grader

for autoregressive language modeling, support ordinary tokenized corpora as sequential environments:

x_1 ... x_n

observation_t = x_1 ... x_t
action_t = distribution/prediction for x_{t+1}
reward_t = log p(x_{t+1} | x_{\le t})

this should permit ordinary next-token training through the SAME external environment interface used for interactive RL tasks.

build adapters/loaders for major open language-model corpora, starting with manageable subsets/shards rather than attempting absurd full-corpus downloads:

* FineWeb / FineWeb-Edu
* Dolma
* C4
* other reputable open LM corpora when useful

FineWeb alone is currently >18T tokens, while Dolma is roughly 3T tokens, so the harness MUST support streaming, sharding, deterministic sampling, train/validation splits, and small reproducible subsets.

also allow arbitrary Hugging Face datasets to be wrapped through a configurable mapping:

```text
dataset row
→ observation symbols
→ target/action symbols
→ reward/grader
```

this should support not only next-token prediction but classification, qa, sequence transformation, code completion, etc.

then integrate REAL interactive environment collections.

PRIORITY 1: integrate every taskset currently exposed by prime intellect's research-environments / verifiers ecosystem that can reasonably be run locally or through its documented runtime.

as of july 22, 2026, prime intellect reports 23 tasksets totaling ~365,000 tasks across software engineering, terminal use, and search. their abstraction separates taskset, harness, and runtime; preserve that separation rather than flattening everything into bespoke adapters.

explicitly investigate and support these prime intellect tasksets:

software engineering:

* swesmith
* openswe
* swerebench_v2
* scaleswe
* swelego
* multiswe
* r2e_gym
* swebench_pro
* swebench_verified
* swebench_multilingual
* senior_swe_bench

terminal:

* tmax
* terminal_lego
* openthoughts_tblite
* terminal_bench_2

search:

* papersearchqa
* wideseek
* s1_deepresearch
* openseeker
* deepdive
* browsecomp
* redsearcher
* browsecomp_plus

prime reports approximately:

* 198k software-engineering tasks
* 28.6k terminal tasks
* 137.6k search tasks

behind one taskset API. use their implementation wherever possible rather than reimplementing datasets or graders.

write an adapter translating their lifecycle into our symbolic environment lifecycle while retaining native rewards, graders, sandboxes, metadata, and termination semantics.

do the same for OTHER established environment zoos.

PRIORITY 2: browser / computer interaction.

integrate BrowserGym as a major environment source. it already exposes multiple web-agent benchmarks through a Gymnasium-style API, including:

* MiniWoB++
* WebArena
* WebArenaVerified
* VisualWebArena
* WorkArena
* AssistantBench
* WebLINX
* OpenApps
* TimeWarp where available

do NOT make separate architectural special cases for each benchmark. build one BrowserGym adapter, then expose benchmark/task identity through configuration.

also investigate OSWorld for real computer/desktop interaction and integrate it through the same symbolic contract. preserve its native action/observation and evaluation semantics underneath the adapter.

PRIORITY 3: standard reinforcement-learning environment zoos.

build a Gymnasium adapter so compatible environments can immediately become symbolic-agent environments.

the Farama ecosystem gives us a large cross-section of useful domains, including:

* MiniGrid
* MiniWorld
* Arcade Learning Environment / Atari
* Gymnasium-Robotics
* MetaWorld
* HighwayEnv
* ViZDoom
* retro environments
* other maintained Gymnasium-compatible environments

Farama explicitly maintains these and lists many additional third-party Gymnasium environments. leverage the ecosystem rather than individually wrapping every environment.

for environments with continuous/native numeric observations or actions, DO NOT quietly feed raw tensors into the symbolic agent and call them symbols.

instead make the boundary explicit:

native observation
→ symbolic encoder
→ symbolic observation

symbolic action
→ native action decoder
→ environment

allow multiple symbolic encoders for the same environment so representation itself can later become an experimental variable.

MiniGrid should be an early target because it is lightweight, configurable, already Gymnasium-compatible, and specifically includes grounded language/instruction tasks.

PRIORITY 4: multi-agent environments.

integrate PettingZoo through a generic adapter.

PettingZoo is effectively a standard API for multi-agent RL and already contains multiple environment families. extend our interaction model from one agent to:

```text
{agent_id → observation symbols}
→ policies
→ {agent_id → action symbols}
→ environment
```

without breaking the single-agent interface.

PRIORITY 5: symbolic/textual embodied environments.

integrate environments that are especially natural fits for this project:

* TextWorld
* ALFWorld
* ScienceWorld

TextWorld is explicitly an extensible sandbox for training/evaluating RL agents in text-generated worlds.

ALFWorld is particularly relevant because it aligns abstract text environments with embodied worlds, allowing high-level policies learned symbolically to correspond to concrete actions.

ScienceWorld provides interactive science tasks requiring grounded experimentation rather than static question answering.

PRIORITY 6: environment standards emerging specifically for agentic RL.

investigate Hugging Face OpenEnv and make our environment abstraction interoperable with it where reasonable.

OpenEnv currently presents a Gymnasium-style `reset` / `step` / `state` protocol for containerized, remotely deployable agent environments and is explicitly intended as a protocol layer for agentic RL. do not unnecessarily compete with that infrastructure; provide an adapter/backend where useful.

the architectural goal should therefore become:

```text
                     ┌─ hf datasets / lm corpora
                     ├─ prime research environments
                     ├─ browsergym
                     ├─ gymnasium / farama
external substrate ──┼─ pettingzoo
                     ├─ osworld
                     ├─ textworld / alfworld / scienceworld
                     ├─ openenv
                     └─ future environment APIs
                              │
                              ▼
                    environment adapter
                              │
                              ▼
                  STANDARD SYMBOLIC ENV API
                              │
                              ▼
                       symbolic agent
```

do not rewrite external environments simply to make the architecture look uniform.

uniformity belongs at the ADAPTER boundary.

retain native:

* reset semantics
* episode boundaries
* rewards
* graders
* action validity rules
* timeouts
* truncation
* sandboxes
* task metadata
* seeds
* dataset splits

and expose normalized versions through our interface.

training must support mixtures of environment families.

i should eventually be able to define something approximately like:

```yaml
training:
  mixture:
    - env: lm/fineweb
      weight: 0.50
    - env: minigrid
      weight: 0.10
    - env: prime/terminal_lego
      weight: 0.10
    - env: prime/scaleswe
      weight: 0.10
    - env: browsergym/miniwob
      weight: 0.10
    - env: textworld
      weight: 0.10
```

and train one compatible agent across the resulting task distribution.

support curriculum schedules as well:

```text
static token prediction
→ simple symbolic prediction
→ short-horizon games
→ text worlds
→ tool use
→ terminal tasks
→ browser tasks
→ software engineering
```

but do NOT assume that curriculum is optimal. make mixture and curriculum policies configurable so we can empirically compare them.

separate TRAINING datasets/environments from HELD-OUT evaluation.

avoid benchmark contamination where possible.

record for every result:

* environment/taskset
* exact version/commit
* split
* seed
* adapter version
* symbolic encoder
* symbolic action decoder
* agent checkpoint
* training data mixture
* reward definition
* episode/token count
* native benchmark metrics

where external benchmarks have canonical scoring, report canonical metrics IN ADDITION TO any normalized reward.

do not replace SWE-bench pass/fail with some homemade universal score, for example. SWE-bench verifies generated patches against repository tests in reproducible environments; preserve that evaluator.

the ultimate experiment is not merely:

"does the framework run?"

it is:

"can one agent architecture learn across a heterogeneous distribution spanning token prediction, static supervised tasks, symbolic worlds, games, search, browsers, terminals, and software engineering when all of them are exposed through a common symbolic interaction protocol?"

build enough integration and training/evaluation infrastructure that we can actually begin answering that question.

also, prime’s list is bigger than i expected: their july 22 post claims the 23 tasksets above and ~365k total tasks, so that alone gives you a pretty absurd downstream distribution once the agent can do nontrivial interaction.

